\section{Introduction}

Modern Transformers \cite{vaswani2017attention,kaplan2020scaling,hoffmann2022chinchilla}
are built around a durable-parameter assumption: large weight tensors are
stored as model state, fetched whenever a layer executes, updated by the
optimizer, and synchronized during distributed training. This assumption is
natural for learning, but not free for systems: in many serving and training
regimes, FFN projections are among the largest repeatedly streamed objects in
the block, and quantization, sparsity, and kernel engineering all preserve the
same lifetime---a parameter tensor remains a durable object the runtime must
move.

This paper asks whether lifetime itself can be part of a kernel interface.  A
generated operand is useful only if its value is created from causal context,
consumed close to the compute unit, and prevented from escaping to persistent
state or ordinary global-memory descriptor storage---a
\emph{no-escape/no-store} contract. The strong claim is inference-scoped:
training additionally needs saved descriptors or deterministic backward
rematerialization. We measure the latter at the FFN level, where
recompute-in-backward cuts the saved-for-backward activation footprint by
$24.5\times$--$71.8\times$ at exact gradient equality. The paper is presented
as a \emph{runtime proposal with measured mechanism evidence and a feasibility
analysis}, not as a completed efficiency method: every claim is labeled
measured, analytic, or unimplemented (Table~\ref{tab:evidencemap}).

We make one architecture the center of the paper.  For each FFN projection
$p$, persistent rank-$r$ factors $U_p,V_p$ surround a small generated core
$M_{p,\mathcal G}$ shared by a causal token group $\mathcal G$:
\begin{equation}
\Pi_{p,\mathcal G}(X_{\mathcal G})=
\bigl((X_{\mathcal G}U_p)M_{p,\mathcal G}\bigr)V_p.
\label{eq:core-intro}
\end{equation}
The canonical implementation uses $M=\operatorname{Diag}(z)$ and emits three
rank vectors $(z_u,z_g,z_d)$.  This is intentionally narrower than generating
complete $(U_t,V_t)$ pairs, but it makes the volatile object small enough to be
a plausible scratch-resident operand.  Full factor synthesis remains a useful
negative boundary because its per-call payload can exceed the dense weights
unless descriptors are reused across large groups.

Three artifacts must not be conflated---the ephemeral-core target (the
proposed method), the historical systems benchmark (isolates low rank, not
lifetime), and the causal software prototype (tests information flow, not
quality or HBM behavior); Table~\ref{tab:artifact-separation},
Appendix~\ref{app:artifacts}, details their evidential roles.

The integrity audit is therefore part of the result rather than a caveat.  The
original controller pooled complete teacher-forced sequences and batches, so
earlier tokens depended on future tokens and unrelated examples; nominal
router seeds did not resample data; and \texttt{torch.where} selected outputs
only after both branches had executed.  We withdraw those quality and traffic
interpretations; the surviving evidence is a persistent-low-rank
microbenchmark, a negative generated/materialized control, corrected
accounting, and software tests for the redesigned causal core.

\begin{definition}[Ephemeral low-rank core]
For a causal token group $\mathcal G$, an ephemeral core is
$M_{p,\mathcal G}=G_\theta(c_{\mathcal G},\ell,p)$ in a declared structured
family $\mathcal M_r$.  It is generated into kernel-local scratch, consumed by
Eq.~\eqref{eq:core-intro}, and is neither returned nor aliased by another
operator, nor committed to persistent state nor stored in a global-memory
descriptor buffer during inference.
\end{definition}

\begin{definition}[Parameter-lifetime class]
A weight-like object has lifetime class $(\tau,\sigma)$, where $\tau$ is its
maximum reuse interval and $\sigma$ is the deepest memory or synchronization
boundary it crosses. Dense FFN weights have long $\tau$ and cross the
model-state boundary; \mwg{} keeps $U_p,V_p,\theta$ durable but assigns each
generated core short $\tau$ and scratch-level $\sigma$.
\end{definition}

\begin{definition}[Eligibility-preserving replacement]
Let $x$ denote a dispatchable whole call and $g(x)\in\{0,1\}$ choose the
ephemeral path.  A replacement is eligibility preserving at call budget
$\routebudget$ if $\E[g(x)]\le\routebudget$, a causal disjoint evaluation is
non-inferior within a predeclared paired margin, and
$\Pr(\Delta(x)>\epsilon\mid g(x)=1)\le\alpha$ for predeclared tail-risk
parameters $(\epsilon,\alpha)$.
\end{definition}

\paragraph{Sampling assumptions for tail-risk certification.}
The tail-risk inequality in Definition~3 is a population statement.  Certifying
it from data requires an explicit sampling model: the selected calls
$\{x: g(x)=1\}$ must be drawn independently from a fixed dispatch
distribution (i.i.d.), and any empirical bound (e.g., a one-sided binomial
upper confidence bound on the failure probability from $n$ selected calls)
inherits only the stated assumptions.  Correlated tokens within a batch,
adaptively selected calls, or distribution shift between calibration and
deployment invalidate the i.i.d.\ premise and require a different
certification route (e.g., concentration under declared mixing, or online
monitoring with sequential tests).  This paper specifies the requirement but
does not certify a numeric tail-risk bound; no empirical failure-probability
claim is made.


\paragraph{Notation.}
We write $d$ for the hidden dimension, $m$ for the FFN intermediate dimension,
$r$ for factor rank, $B$ for batch size, $S$ for sequence length, and $T=BS$
for tokens per FFN call.  A core is shared by a causal group of $q$ tokens.
$\Gamma(r)=r(d+m)/(dm)$ is the persistent-factor-to-dense element ratio; it is
not a device-traffic measurement.  For the proposed router, excess loss is
$\Delta(x)=\loss_{\mathrm{mwg}}(x)-\loss_{\mathrm{dense}}(x)$: smaller values
favor patching, and an eligibility bound patches only when
$\widehat\Delta(x)\le\Delta_{\max}$. Payload is $\traffic$,
loss is $\loss$, and expectation is $\E$.

\paragraph{Contributions.}
\begin{itemize}
  \item An inference-side no-escape/no-store operator for ephemeral low-rank
  cores, separated from persistent factors and training rematerialization;
  three token-local diagonal cores replace six non-identifiable
  paired scales, removing one explicitly described redundancy;
  causality, algebra, and checkpoint behavior are verified in
  software tests.
  \item Call-level payload and reuse break-even conditions, whole-call versus
  token-fraction dispatch, and a repaired approximation result whose retained
  Frobenius energy matches its matrix norm.
  \item The fused no-store arm as an audited, hash-bound Ascend-C operator,
  and the predeclared four-arm experiment run on device: algebraic
  equivalence at the fp16 rounding floor, core DRAM traffic verified
  output-only, and \emph{parity} with the materialized arm at L2-resident
  geometry after host fast-pathing (beyond-L2 deficit isolated to a
  single-core epilogue; the measured v4 multi-core rewrite is
  bit-identical to v1 and $1.13\times$--$3.03\times$ faster across
  $T{=}128$--$4096$; Section~\ref{sec:systems}).  Separately, on the training
  side (where the core is recomputed, not stored), the low-rank
  parameterization yields an optimizer-inclusive
  training step
  synchronized across four \texttt{DistributedDataParallel} workers
  that runs $3.37\times$--$4.14\times$ faster than dense---a rank effect, not
  a no-store effect.
  \item Withdrawal of invalid quality/routing evidence, plus a scoped
  persistent-low-rank measurement and an unfavorable conditioning control,
  for context.
\end{itemize}

